\documentclass[12pt,letterpaper]{article}
\usepackage[utf8]{inputenc}
\usepackage[spanish]{babel}
\usepackage{amsmath}
\usepackage{amsfonts}
\usepackage{amssymb}
\usepackage[margin=2.5cm]{geometry}

\begin{document}

\begin{center}
    \Large \textbf{Evaluación de Examen 2: Unidad II} \\
    \large Física para Ciencias de la Salud (005-1713) \\
    \vspace{0.5cm}
    \textbf{Estudiante:} Ana Alcalá \quad \textbf{C.I.:} 30.749.815
    \vspace{0.3cm} \\
    \large \textbf{Calificación Final: 10/10 puntos}
\end{center}

\vspace{0.5cm}

\section*{Análisis de Resultados}

\subsection*{1. Cinemática (Aceleración media)}
\textbf{Procedimiento y resultado:} La estudiante identifica correctamente los datos iniciales ($v_0 = 0 \text{ m/s}$, $v_f = 0,6 \text{ m/s}$, $t = 0,15 \text{ s}$) y aplica la fórmula de la aceleración media. El cálculo aritmético $0,6 / 0,15$ arroja el resultado correcto de $4 \text{ m/s}^2$. \\
\textbf{Veredicto:} \textbf{Correcto} (2/2 pts).

\subsection*{2. Dinámica (Leyes de Newton)}
\textbf{Procedimiento y resultado:} Plantea de forma impecable la Segunda Ley de Newton ($F = m \cdot a$). Realiza un despeje analítico correcto para la aceleración ($a = \frac{F_{\text{net}}}{m}$) y sustituye los valores numéricos para obtener la aceleración exacta de $6 \text{ m/s}^2$. \\
\textbf{Veredicto:} \textbf{Correcto} (2/2 pts).

\subsection*{3. Trabajo Mecánico}
\textbf{Procedimiento y resultado:} Demuestra un excelente dominio teórico al identificar que la fuerza y el desplazamiento son paralelos, por lo que el ángulo es $0^\circ$. Sustituye correctamente $\cos(0^\circ) = 1$ en la fórmula general del trabajo ($W = F \cdot d$). Su cálculo final de $320 \text{ J}$ es perfectamente exacto. \\
\textbf{Veredicto:} \textbf{Correcto} (2/2 pts).

\subsection*{4. Energía Potencial}
\textbf{Procedimiento y resultado:} Extrae apropiadamente los datos de masa ($1,2 \text{ kg}$) y altura ($1,5 \text{ m}$), utilizando de manera adecuada el valor de la gravedad estándar de $9,8 \text{ m/s}^2$. Efectúa la multiplicación de los tres términos de la fórmula de energía potencial gravitatoria, obteniendo el resultado correcto de $17,64 \text{ J}$. \\
\textbf{Veredicto:} \textbf{Correcto} (2/2 pts).

\subsection*{5. Potencia Mecánica}
\textbf{Procedimiento y resultado:} Aplica directamente la fórmula de la potencia mecánica al relacionar el trabajo efectuado ($75 \text{ J}$) entre el intervalo de tiempo ($60 \text{ s}$). La operación se desarrolla con éxito, brindando la cifra final correcta de $1,25 \text{ W}$. \\
\textbf{Veredicto:} \textbf{Correcto} (2/2 pts).

\vspace{0.5cm}
\section*{Comentarios Generales y Sugerencias}
¡Felicidades por alcanzar la calificación máxima en la Unidad II! La estudiante evidencia un nivel sobresaliente y gran razonamiento analítico.
\begin{itemize}
    \item El examen muestra un nivel de orden ejemplar; en todo momento se utilizaron correctamente los despejes analíticos antes de las sustituciones numéricas.
    \item Se valora positivamente el detalle de incluir siempre las magnitudes entre paréntesis (como los Joules o los Watts) al declarar el resultado final, un excelente hábito propio de las ciencias de la salud.
    \item Se invita a la estudiante a mantener el mismo nivel de dedicación para las siguientes unidades.
\end{itemize}

\end{document}